\begin{figure}[tbp]
\centering
\begin{figurealt}{CMake receives a default renderer and optionally a list of renderers to compile. A generated registry maps public identities to family factories. Program preference, environment and default resolve to a selected identity; the first graphics device latches the active family. A family may then call a native API or library. Five identities map to the single easygl family.}
\resizebox{\textwidth}{!}{%
\begin{tikzpicture}[node distance=9mm and 13mm]
  \node[cna public] (options) {default + compiled set};
  \node[cna internal, right=of options] (registry) {generated identity /\\factory registry};
  \node[cna public, right=of registry] (policy) {explicit preference $>$\\environment $>$ default};
  \node[cna internal, right=of policy] (family) {first device latches\\active family};
  \node[cna public, right=of family] (api) {native API / library};
  \node[cna public, below=of registry] (five) {OPENGLES2, OPENGLES3,\\OPENGL33, WEBGL1, WEBGL2};
  \node[cna internal, below=of family] (easy) {one \texttt{easygl} family\\runtime profile};
  \draw[cna flow] (options) -- (registry);
  \draw[cna flow] (registry) -- (policy);
  \draw[cna flow] (policy) -- (family);
  \draw[cna flow] (family) -- (api);
  \draw[cna flow] (five) -- (easy);
  \draw[cna flow] (easy) -- (api);
\end{tikzpicture}}
\end{figurealt}
\caption{Compiled set, default, runtime selection, active family, and native API are distinct.
The five-to-one EasyGL mapping explains the \RendererIdentityCount{}-identity /
\RendererFamilyCount{}-family count.}
\label{fig:renderer-identity-family-api}
\end{figure}
